%% --- 6.5 Review Scores ---

\subsection{Review Scores}
\label{sec:exp:reviews}

We evaluate perceived paper quality using ScholarPeer~\citep{goyal2026scholarpeer}, an automated peer review system backed by \texttt{gemini-3.1-pro-preview} and rich literature search support.

\begin{table}[h]
\centering
\small
\caption{ScholarPeer review rating scores (1--4 scales except Overall (1-10 scale)) and accept decisions. \emph{Average}: mean across 15~papers per system. \emph{Best-of-3}: strongest seed per task.}
\label{tab:reviews}
\begin{tabular}{@{}lcccccc@{}}
\toprule
\textbf{System} & \textbf{Soundness} & \textbf{Originality} & \textbf{Quality} & \textbf{Clarity} & \textbf{Overall} & \textbf{\#Accept} \\
\midrule
\multicolumn{7}{l}{\textit{Average (3 seeds $\times$ 5 tasks)}} \\
Sakana AI-Scientist v2    & 1.5 & 1.9 & 1.5 & 3.1 & 2.5 & 0/15 \\
AutoResearchClaw            & 1.1 & 2.3 & 1.1 & 2.5 & 1.9 & 0/15 \\
DeepScientist             & 1.7 & 1.7 & 1.6 & 3.1 & 2.5 & 1/15 \\
AI-Researcher            & 1.9 & 2.4 & 1.9 & 3.1 & 3.4 & 2/15 \\
\sys{}         & \textbf{2.3} & \textbf{2.5} & \textbf{2.3} & 3.0 & \textbf{4.5} & \textbf{6/15} \\
\midrule
\multicolumn{7}{l}{\textit{Best-of-3 (strongest seed per task)}} \\
Sakana AI-Scientist v2    & 1.6 & 2.0 & 1.6 & 3.2 & 3.4 & 0/5 \\
AutoResearchClaw            & 1.2 & 2.0 & 1.2 & 3.0 & 3.0 & 0/5 \\
DeepScientist             & 2.2 & 2.0 & 2.2 & 3.6 & 3.6 & 1/5 \\
AI-Researcher            & 2.0 & 2.4 & 2.0 & 3.2 & 4.0 & 1/5 \\
\sys{}         & \textbf{2.8} & \textbf{3.0} & \textbf{2.8} & \textbf{3.6} & \textbf{6.6} & \textbf{4/5} \\
\bottomrule
\end{tabular}
\end{table}



\paragraph{Verifiable papers are deemed better by automatic reviewers.}
\sys{} achieves a 40\% accept rate (6/15), tripling the best baseline (AIR: 13\%), and best-of-3 selection reaches 6.6 overall rating score and 4/5~tasks accepted.
This gap is not driven by better algorithms---solver scores cluster tightly across systems (Table~\ref{tab:solver})---but by what happens \emph{after} the solver finishes.
The Claim Verifier prevents the most damaging failure mode we observe in rejected papers: claims that contradict the paper's own data (e.g., ``sub-millisecond latency'' when the results table reports 7.9\,ms).
% Notably, \coeaudit{} and review scores measure orthogonal properties: across all 75~papers, Clarity averages 1--2 points above Soundness, confirming that narrative coherence is independent of evidence integrity.

\paragraph{The paper quality is bottlenecked by research soundness, not writing capability.}
Across all systems, Clarity is consistently the highest-scoring dimension (2.5--3.1) while Soundness is the lowest (1.1--2.3): these papers read well but do not withstand methodological scrutiny.
The reviewer's two most frequent complaints are missing comparisons against published baselines and proxy-only evaluation without end-to-end system measurements. While \sys{}'s Problem Investigator retrieves related work and identifies candidate baselines, the resulting comparisons do not yet meet the depth that ScholarPeer expects (e.g., re-implementing a SOTA method and reporting head-to-head numbers).
\sys{} also exhibits high seed variance (e.g., EPLB scores 1, 3, 8 across three seeds on the same task): rejected runs are those where the paper writer generates claims that the Claim Verifier's current coverage does not fully catch---for example, exaggerated qualitative framing (``near-optimal'') rather than numerically falsifiable statements.
Accepted runs make calibrated claims from the same underlying data, suggesting that extending verification coverage to qualitative claims would reduce this variance.



